% header.tex: pandoc header-includes for report.
% Polishes tables, links, code, headings, and margins using only
% widely-available LaTeX packages (texlive-latex-base/recommended).
% Verified against: booktabs, longtable, tabularx, array, colortbl,
% xcolor, hyperref, caption, fancyhdr, fancyvrb.

% ---- Tables ----
\usepackage{booktabs}      % \toprule \midrule \bottomrule
\usepackage{longtable}     % multi-page tables (pandoc default)
\usepackage{tabularx}      % X column type with auto-wrap
\usepackage{array}         % >{...} column prefix, \arraybackslash
\usepackage{colortbl}      % row coloring
\setlength{\tabcolsep}{3pt}     % was 6pt; column widths assume a small inter-column budget
\renewcommand{\arraystretch}{1.15}
\setlength{\arrayrulewidth}{0.3pt}

% Two-line subtitle helper: each line stays centered in the title block.
% No vspace between lines (so they read as one breath); \par alone for the
% break opportunity.
\newcommand{\subtwoline}[2]{#1\par\nobreak #2}

% Pre-empt pandoc's default \subtitle by defining it ourselves first;
% pandoc's later \providecommand will then no-op. The \vspace gives the
% title room to breathe before the subtitle block starts.
\makeatletter
\newcommand{\subtitle}[1]{%
  \apptocmd{\@title}{\par\vspace{0.9em}{\large #1\par}}{}{}%
}
\makeatother

% TeX Gyre Pagella's right-single-quote (U+2019) glyph is visually heavy;
% pandoc's smart-quote conversion produces \textquoteright from every
% apostrophe in body text, which then renders as that heavy glyph. Remap
% to the typewriter straight apostrophe \textquotesingle for a leaner
% look. Applies uniformly to all body-text apostrophes.
\renewcommand{\textquoteright}{\textquotesingle}

% Targeted footnote spacing: \thanks footnotes (\fnsymbol) stay compact;
% only the first numeric body footnote on a page gets extra leading space,
% so the † → 1 transition has visual separation while * → † stays tight.
% Trick: wrap \@makefntext to prepend a vskip when the marker is numeric
% AND the previous footnote was a symbolic one. We track that via a flag.
\makeatletter
\newif\if@prevfnsymbolic
\@prevfnsymbolictrue   % start true: first body footnote on title page benefits

% Save the original \thanks (article-class defined macro that mutates
% the title's thanks-list and bumps the footnote counter).
\let\@hdr@orig@thanks\thanks
\renewcommand{\thanks}[1]{%
  \@hdr@orig@thanks{#1}%
  \global\@prevfnsymbolictrue}

\let\@hdr@orig@makefntext\@makefntext
\renewcommand{\@makefntext}[1]{%
  \ifnum\value{footnote}>0
    \if@prevfnsymbolic
      \vspace{0.5em}%
      \global\@prevfnsymbolicfalse
    \fi
  \fi
  \@hdr@orig@makefntext{#1}}
\makeatother
% Note: \arrayrulecolor{ruleline!35} is set further down, after \definecolor.

% ---- Color theme ----
\usepackage{xcolor}
\definecolor{accent}{HTML}{1F6FB5}     % steel-blue
\definecolor{accent2}{HTML}{B85C00}    % burnt-orange (warnings)
\definecolor{rowalt}{HTML}{F4F7FA}     % subtle alternating row
\definecolor{ruleline}{HTML}{4A6076}
\arrayrulecolor{ruleline!35}    % faint inter-row rules in longtables

% ---- microtype: protrusion, expansion, character-level tracking ----
% Subtle typographic polish: small glyph protrusion into the right margin,
% imperceptible font expansion to improve justification, and smarter
% character tracking. Works out of the box with lualatex + fontspec.
\usepackage{microtype}

% ---- Widow / orphan control ----
% widows-and-orphans (Mittelbach) does two things: (1) `[prevent-all]`
% sets \clubpenalty, \widowpenalty, \displaywidowpenalty, \brokenpenalty
% all to 10 000, forbidding the page-breaker from creating any of them;
% (2) reports any that nevertheless slip through in the .log so we can
% deal with them by hand. The package is bundled locally; see Makefile.
\usepackage[prevent-all]{widows-and-orphans}
\definecolor{codebg}{HTML}{F4F4EE}
\definecolor{codeborder}{HTML}{D5D9DC}

% ---- Hyperlinks: colored, not boxed ----
\usepackage[colorlinks=true,
            linkcolor=accent,
            urlcolor=accent,
            citecolor=accent,
            allbordercolors={1 1 1}]{hyperref}

% ---- Heading styles via plain \makeatletter overrides (no titlesec) ----
\makeatletter
\renewcommand\section{\@startsection{section}{1}{\z@}%
  {-3.5ex plus -1ex minus -.2ex}%
  {2.0ex plus .2ex}%
  {\normalfont\Large\bfseries\color{accent}}}
\renewcommand\subsection{\@startsection{subsection}{2}{\z@}%
  {-2.5ex plus -1ex minus -.2ex}%
  {1.2ex plus .2ex}%
  {\normalfont\large\bfseries\color{accent}}}
\renewcommand\subsubsection{\@startsection{subsubsection}{3}{\z@}%
  {-2.0ex plus -.7ex minus -.2ex}%
  {.8ex plus .2ex}%
  {\normalfont\normalsize\bfseries\color{accent!75!black}}}
\makeatother

% ---- Code (verbatim and inline) ----
% fvextra extends fancyvrb with breaklines support. Bundled locally
% (fvextra.sty, upquote.sty) so the build works without modifying the
% texlive install. fvextra requires upquote.
\usepackage{upquote}
\usepackage{fvextra}
% Verbatim blocks: framed, subtle background, break long lines at
% whitespace AND at characters when needed. Applies to pandoc's
% Highlighting environment via inherited \fvset defaults.
\fvset{frame=single,
       rulecolor=\color{codeborder},
       framesep=4pt,
       fontsize=\small,
       breaklines=true,
       breakanywhere=true,
       breakindent=2em,
       samepage=true}        % no page-break inside a Verbatim block

% Pandoc defines `Highlighting` via \DefineVerbatimEnvironment WITHOUT
% picking up our \fvset{samepage=true} global default (it overrides
% \FV@KeyValues with just commandchars). Re-bind it with samepage in the
% per-env option list so \interlinepenalty=10000 is in scope inside the
% env.
\usepackage{etoolbox}
\AtBeginDocument{%
  \DefineVerbatimEnvironment{Highlighting}{Verbatim}%
    {commandchars=\\\{\},samepage=true}%
}

% fvextra @Break (active under breaklines=true) emits an anti-aliasing
% overlap glue (\vspace{-\FV@backgroundcolorboxoverlap}, default 0.25pt)
% AFTER each line's hbox via \FV@bgcoloroverlap (fvextra.sty:3698). That
% glue node sits UPSTREAM of fancyvrb's \penalty\interlinepenalty (which
% goes before the next box). TeX's break-finding sees the glue as a
% feasible candidate at p=0, ignoring the downstream penalty. samepage
% therefore has no teeth on syntax-highlighted blocks.
%
% Why a naive `\nobreak\vspace{-X}` isn't enough:
%   fancyvrb's @i (the first-line dispatcher) wraps line-processing in an
%   extra `\hbox{}`, so the trailing \FV@bgcoloroverlap runs in restricted
%   HMODE. In hmode, `\vspace{-X}` expands to `\vadjust{\vskip-X\vskip 0pt}`
%   -- the \vskip nodes get queued for emission into the OUTER vlist AFTER
%   the line's hbox is shipped. A preceding `\nobreak` in hmode emits a
%   penalty INSIDE the hbox -- which the page-breaker can't see (penalties
%   inside hboxes don't affect page-breaking). Net result: the \vskip
%   lands in the outer vlist unprotected, and the page-breaker still sees
%   a p=0 candidate immediately after the first verbatim line.
%
% Fix: build the protective penalty + glue as one atomic vlist unit. In
% vmode, emit \penalty\@M then \vskip directly. In hmode (the @i case),
% wrap BOTH inside a single \vadjust so they enter the outer vlist
% together. \vspace's hmode/vmode duality is the trap; bypassing \vspace
% removes it. The overlap remains visually (PDF anti-aliasing trick to
% hide seams between line backgrounds) but the page-breaker never sees
% it as a feasible candidate.
\makeatletter
\def\FV@bgcoloroverlap{%
  \ifx\FV@backgroundcolorboxoverlap\relax
  \else
    \ifvmode
      \penalty\@M
      \vskip-\FV@backgroundcolorboxoverlap\relax
    \else
      \vadjust{\penalty\@M\vskip-\FV@backgroundcolorboxoverlap\relax}%
    \fi
  \fi}
\makeatother

% ---- Captions ----
\usepackage[font=small,labelfont=bf,labelsep=period]{caption}

% ---- Breakable \texttt for table cells ----
% Long identifiers like applyCorrectionsPolarAndConvertToCartesian are
% atomic tokens that won't break to fit narrow columns. Workaround: a
% pandoc Lua filter (filter-codebreaks.lua) rewrites \texttt{...} content
% to insert \allowbreak{} at separators (::, _, ., -, /, () and camelCase
% boundaries. Identifiers without any natural break may still bleed a
% few points; that's the cost of all-lowercase compound names.

% ---- Margins ----
\usepackage[margin=0.9in,top=0.85in,bottom=1.0in]{geometry}

% ---- Header / footer ----
\usepackage{fancyhdr}
\pagestyle{fancy}
\fancyhf{}
% Running header text. Set \runninghead in your document (or via pandoc
% header-includes) to put a label top-left; defaults to empty.
\providecommand{\runninghead}{}
\fancyhead[L]{\small\color{ruleline}\runninghead}
\fancyhead[R]{\small\color{ruleline}\thepage}
\fancyfoot[C]{}
\renewcommand{\headrulewidth}{0.4pt}
\renewcommand{\footrulewidth}{0pt}
\setlength{\headheight}{14pt}
\setlength{\headsep}{18pt}

% ---- Pandoc-friendly tightlist ----
% Pandoc emits \tightlist for any list whose items lack blank lines between
% them. Default zeros itemsep + parskip; we leave a small itemsep so even
% "tight" lists with multi-sentence items breathe a little.
\providecommand{\tightlist}{\setlength{\itemsep}{0.35em}\setlength{\parskip}{0pt}}

% ---- Diff syntax highlighting: green for +, red for - ----
% Pandoc's skylighting maps ```diff added/removed lines to \VariableTok and
% \StringTok respectively, which the default pygments style renders as two
% nearly identical blues (RGB 0.10/0.09/0.49 vs 0.25/0.44/0.63), invisible
% as a diff. Override with conventional green/red so + and - read at a
% glance. C++ string literals (also \StringTok) become red, which matches
% common IDE themes; variable highlighting in other languages (e.g. shell)
% becomes green. Pandoc emits the \newcommand{...} definitions ahead of
% the include-in-header content, so \renewcommand here lands cleanly.
\definecolor{diffaddcolor}{HTML}{007020}
\definecolor{diffdelcolor}{HTML}{B22222}
\renewcommand{\VariableTok}[1]{\textcolor{diffaddcolor}{#1}}
\renewcommand{\StringTok}[1]{\textcolor{diffdelcolor}{#1}}

% ---- longtable spacing ----
\setlength{\LTpre}{6pt}
\setlength{\LTpost}{6pt}

% ---- Graphics extension priority: prefer PDF over PNG ----
% (allows extensionless image refs in markdown, or works alongside
%  the report's explicit `.pdf` refs.)
\DeclareGraphicsExtensions{.pdf,.png,.jpg}

% ---- Global graphics: cap to text area, keep aspect ratio ----
% Matplotlib charts are typically saved at a fixed inch size that's
% wider than this report's text column; without this they overflow
% the right margin.
\setkeys{Gin}{width=\linewidth,keepaspectratio}

% ---- Pandoc passes `figure`-floats through; let them break across pages ----
\usepackage{float}
\floatplacement{figure}{H}

% ---- A "softnote" for cross-reference paragraphs ----
\providecommand{\softnote}[1]{\par\medskip\noindent{\itshape\color{ruleline}#1}\par\medskip}
